%% kapa: exam=1 round=35 year=2024 period=2 form=A subject=accounting subject_ko=회계학 seq=29 number=69
%% source: Q-Net 감정평가사 기출문제 (cst00309) artlSeq=5213229 file=raw/1008/5213229/2교시 문제지 원본.pdf
%% conversion: kordoc 4.24.51 -> tools/gen1.py (auto); review=proofread-sonnet; flags=none
\begin{question}{69}
(주)감평은 20×1년 4월 1일 업무용 기계장치를 취득(취득원가 ￦61,000, 내용연수 5년, 잔존가치 ￦1,000)하여 정액법으로 감가상각하였다. (주)감평은 20×2년 10월 1일 동 기계장치의 감가상각방법을 연수합계법으로 변경하고 남은 내용연수도 3년으로 재추정하였으며, 잔존가치는 변경하지 않았다. 20×2년도 포괄손익계산서에 인식할 기계장치의 감가상각비는? (단, 동 기계장치는 원가모형을 적용하며, 감가상각은 월할계산한다.)
\begin{choices}
  \item ￦5,250
  \item ￦9,150
  \item ￦12,200
  \item ￦13,250
  \item ￦14,250
\end{choices}
\end{question}
\answerall
\answernote{정답표에 1,2,3,4,5 모두 정답으로 표기(전원 정답)}
